% Rozdział współdzielony tematycznie z geometrią różniczkową; wersja audytowana 2026-11-19.
% Samodzielny rozdział do włączenia przez \input{rozdzial_dzety.tex}.
% Wymaga: amsmath, amssymb, amsthm, mathtools, tikz, hyperref.
% Środowiska: definicja, twierdzenie, lemat, wniosek, uwaga, przyklad, proof.
\chapter{Dżety: od geometrii różniczkowej do teorii miary}
\label{ch:dzety}

\section{Motywacja: ile informacji niesie punkt?}
W rozdziale~\ref{ch:area-coarea-28} wyprowadziliśmy twierdzenie
Rademachera i wzory miarowo-geometryczne. Teraz interpretujemy
różniczkowalność prawie wszędzie językiem dżetów, zachowując też
pełne ujęcie geometryczne, formy kontaktowe i twierdzenia Whitneya.
Wartość $f(p)$ opisuje jedynie punkt obrazu, pochodna $Df(p)$ ---
przybliżenie liniowe, a kolejne pochodne --- coraz dokładniejsze
przybliżenia lokalne. \emph{Dżet rzędu $k$} jest sposobem zebrania
tej informacji w jeden obiekt. Rozróżniamy trzy sytuacje:
\begin{enumerate}
\item w geometrii gładkiej dżet jest klasą równoważności
kiełków odwzorowań lub przekrojów;
\item w równaniach różniczkowych współrzędne dżetowe
traktuje się jako formalnie niezależne, lecz rozwiązania
muszą spełniać warunki holonomiczności;
\item w teorii miary dżet może istnieć tylko prawie wszędzie,
a zwykłą granicę zastępuje się \emph{granicą przybliżoną}.
\end{enumerate}
Te trzy punkty widzenia spotykają się w twierdzeniu Whitneya:
\emph{kiedy formalnie zadane pochodne dają się zrealizować przez
prawdziwą funkcję gładką?} Wersja Whitneya--Luzina mówi zaś,
że wiele funkcji mierzalnych staje się gładkimi po pominięciu
zbioru o dowolnie małej mierze.

\section{Kiełki, klasy równoważności i wielomiany Taylora}
\begin{definicja}[Kiełek odwzorowania]
\label{def:dzety-kielek}
Niech $M,N$ będą gładkimi rozmaitościami i $p\in M$.
Dwa gładkie odwzorowania $f:U\to N$ i $g:V\to N$,
gdzie $p\in U\cap V$, mają ten sam \emph{kiełek w $p$},
jeżeli pokrywają się na pewnym otoczeniu $p$ zawartym w $U\cap V$.
Zbiór kiełków oznaczamy przez $\mathcal C^\infty_p(M,N)$.
\end{definicja}

\begin{definicja}[$k$-dżet odwzorowania]
\label{def:dzety-klasa}
Niech $M,N$ będą gładkimi rozmaitościami, $p\in M$,
a $k\in\mathbb N_0$. Na zbiorze kiełków
$\mathcal C^\infty_p(M,N)$ definiujemy relację:
kiełki o różnych wartościach w $p$ nie są równoważne.
Jeżeli $f(p)=g(p)=q\in N$, piszemy $f\sim_{p,k}g$, gdy
po wybraniu mapy $x$ wokół $p$ i mapy $y$ wokół $q$ zachodzi
\[
 \partial^\alpha (y\circ f\circ x^{-1})(x(p))
 =\partial^\alpha (y\circ g\circ x^{-1})(x(p)),
 \qquad |\alpha|\le k.
\]
Dziedziny przedstawicieli ograniczamy tak, aby obrazy
leżały w mapie celu. Zgodność pochodnych nie zależy
od przedstawicieli kiełków, gdyż są oni równi w pewnym
otoczeniu $p$. Niezależność od map wykażemy poniżej.
Klasę równoważności oznaczamy przez $j_p^k f$ i nazywamy
\emph{$k$-dżetem $f$ w $p$}. Dla $k=0$ jest to po prostu para
$(p,f(p))$. Wartości pochodnych liczymy składowo.
\end{definicja}

\begin{lemat}[Niezależność od map]
\label{lem:dzety-mapy}
Relacja w Definicji~\ref{def:dzety-klasa} nie zależy od wyboru
map źródła i celu.
\end{lemat}
\begin{proof}
Niech $\varphi,\psi$ będą zmianami współrzędnych.
Dla odwzorowania współrzędnego $F$ nowy zapis ma postać
$\psi\circ F\circ\varphi^{-1}$. Reguła łańcuchowa mówi,
że jego pochodna rzędu $r\le k$ jest skończoną sumą wyrażeń
zawierających pochodne $F$ rzędów od $1$ do $r$, pochodne
$\psi,\varphi^{-1}$ i wartości funkcji w rozpatrywanym punkcie.
Jeżeli odpowiednie dane dla $F$ i $G$ są równe do rzędu $k$,
to równe są także wszystkie nowe pochodne. Zmieniając mapy
w przeciwnym kierunku, otrzymujemy równoważność.
\end{proof}

\begin{przyklad}[Różne funkcje, ten sam dżet]
\label{ex:dzety-taylor}
Dla $f(x)=\sin x$ i $g(x)=x$ mamy $j_0^2f=j_0^2g$,
ale $j_0^3f\ne j_0^3g$. Natomiast $f(x)=x+x^2$
i $g(x)=x-x^2$ mają identyczne $1$-dżety w zerze,
ale różne $2$-dżety.
\end{przyklad}

\begin{uwaga}[Nie należy utożsamiać drugiej pochodnej z tensorem]
\label{rem:dzety-hess}
Pierwsza pochodna funkcji skalarnych definiuje kanoniczny
kowektor $\mathrm df_p\in T_p^*M$. Dla drugiej pochodnej jest inaczej.
Współrzędna $x$ na $\mathbb R$ i współrzędna $y=x+x^2$ wokół zera
dałyby dla funkcji $f(x)=x$ odpowiednio $f''(0)=0$
oraz $\frac{\mathrm d^2 f}{\mathrm dy^2}(0)=-2$.
Oznacza to, że $2$-dżet jest obiektem niezmienniczym,
ale same współrzędne drugich pochodnych nie są tensorami.
Tensor Hessego powstaje kanonicznie dopiero po określeniu
koneksji (lub w punkcie krytycznym funkcji).
\end{uwaga}

\begin{definicja}[Algebra dżetów funkcji skalarnych]
\label{def:dzety-algebra}
Niech $\mathcal E_p$ oznacza pierścień kiełków gładkich
funkcji rzeczywistych na $M$ w punkcie $p$ i
$\mathfrak m_p=\{[f]:f(p)=0\}$ jego ideał maksymalny.
Definiujemy $\mathcal J_p^k=\mathcal E_p/\mathfrak m_p^{k+1}$.
\end{definicja}

\begin{twierdzenie}[Algebra i wielomiany Taylora]
\label{thm:dzety-algebra}
Dla $M=\mathbb R^n$ i $p=0$ odwzorowanie
\[
 [f]\longmapsto \sum_{|\alpha|\le k}
 \frac{\partial^\alpha f(0)}{\alpha!}X^\alpha
\]
indukuje izomorfizm algebr
$\mathcal J_0^k\cong\mathbb R[X_1,\ldots,X_n]/(X_1,\ldots,X_n)^{k+1}$.
W szczególności $\dim\mathcal J_0^k=\binom{n+k}{k}$.
\end{twierdzenie}
\begin{proof}
Wielomian Taylora rzędu $k$ iloczynu jest iloczynem
wielomianów Taylora po odrzuceniu wyrazów stopnia większego niż $k$,
co wynika z reguły Leibniza. Jądro odwzorowania składa się
z kiełków o pochodnych rzędów $0,\ldots,k$ równych zeru.
Wybieramy przedstawiciela $f$ w kuli wokół zera.
Dla ustalonego $x$ stosujemy jednowymiarowy wzór
Taylora z resztą całkową do $t\mapsto f(tx)$.
Ponieważ pochodne do rzędu $k$ w zerze znikają,
\[
 f(x)=\sum_{|\alpha|=k+1}x^\alpha
 \frac{k+1}{\alpha!}\int_0^1(1-t)^k
 \partial^\alpha f(tx)\,\mathrm d\lambda_1(t).
\]
Współczynniki przy $x^\alpha$ są gładkie:
na każdej mniejszej domkniętej kuli wszystkie
pochodne całkowanej funkcji względem $x$ są
ciągłe i ograniczone, więc wolno różniczkować
pod całką. Każdy $x^\alpha$ jest iloczynem $k+1$
kiełków z $\mathfrak m_0$, co daje inkluzję
jądra w $\mathfrak m_0^{k+1}$.
Odwrotnie, w każdej pochodnej rzędu co najwyżej $k$
iloczynu $k+1$ kiełków z $\mathfrak m_0$ reguła
Leibniza pozostawia co najmniej jeden czynnik
niezróżniczkowany, który znika w zerze.
To samo dotyczy skończonych sum takich iloczynów.
Zatem jądrem jest dokładnie $\mathfrak m_0^{k+1}$.
Każdy wielomian stopnia co najwyżej $k$ sam jest gładki,
więc homomorfizm jest surjektywny. Liczba jednomianów
o $|\alpha|\le k$ wynosi $\binom{n+k}{k}$.
\end{proof}

\begin{przyklad}[Liczby dualne]
\label{ex:dzety-dualne}
Dla $n=k=1$ dostajemy
$\mathcal J_0^1\cong\mathbb R[\varepsilon]/(\varepsilon^2)$.
Dżet funkcji $f$ ma zapis $f(0)+f'(0)\varepsilon$;
reguła mnożenia automatycznie daje wzór na pochodną iloczynu.
To algebraiczne źródło arytmetyki liczb dualnych używanej
w automatycznym różniczkowaniu.
\end{przyklad}

\section{Wiązki dżetów i ich naturalne odwzorowania}
\begin{definicja}[Wiązka dżetów odwzorowań]
\label{def:dzety-bundle}
Dla gładkich rozmaitości $M,N$ oznaczamy przez $J^k(M,N)$
zbiór wszystkich $k$-dżetów lokalnych gładkich odwzorowań
$M\to N$. Naturalne odwzorowania to
$\pi_0^k(j_p^k f)=(p,f(p))$ oraz dla $0\le r\le k$
$\pi_r^k(j_p^k f)=j_p^r f$.
Dla ustalonego $f$ odwzorowanie $j^kf:p\mapsto j_p^kf$
nazywamy \emph{przedłużeniem dżetowym}.
\end{definicja}

W lokalnych współrzędnych na $M$ i $N$, o wymiarach $n,m$,
przestrzeń $J^k(M,N)$ ma współrzędne $(x^i,u^a,u^a_\alpha)$,
$1\le|\alpha|\le k$, gdzie $u^a_\alpha$ odpowiada pochodnej
$\partial^\alpha u^a$. Jej wymiar wynosi
$n+m\binom{n+k}{k}$. Każdy zestaw współrzędnych nad ustalonymi punktami
map jest realizowany przez wielomian Taylora
o zadanych współczynnikach, po ograniczeniu
otoczenia punktu źródłowego tak, aby obraz leżał
w mapie celu. Współrzędne są więc bijekcjami
na otwarte podzbiory przestrzeni euklidesowej.
Zmiany współrzędnych są gładkie na mocy reguły
łańcuchowej i mają gładkie odwrotności.
Topologia określona przez te mapy jest Hausdorffa:
dżety o różnych punktach bazowych w $M\times N$
rozdzielamy przez przeciwobrazy rozłącznych otoczeń,
a różne dżety nad tym samym punktem rozdzielamy
w jednej wspólnej mapie dżetowej.
Przeliczalne atlasy $M,N$ dają przeliczalną rodzinę
map dżetowych, każdą z przeliczalną bazą; stąd drugi
aksjomat przeliczalności. Otrzymujemy rozmaitość gładką.

\begin{twierdzenie}[Pierwsze dżety i struktura afiniczna]
\label{thm:dzety-affine}
Dla dowolnych $(p,q)\in M\times N$ włókno
$(\pi_0^1)^{-1}(p,q)$ utożsamia się kanonicznie z
$\operatorname{Hom}(T_pM,T_qN)$. Dla $k\ge2$ włókna
$\pi_{k-1}^k:J^k(M,N)\to J^{k-1}(M,N)$ są przestrzeniami
afinicznymi modelowanymi na
$S^kT_p^*M\otimes T_qN$, gdzie $(p,q)$ jest
wartością rozpatrywanego $(k-1)$-dżetu.
\end{twierdzenie}
\begin{proof}
Wartość i różniczka $f$ w $p$ jednoznacznie wyznaczają
$1$-dżet, a dowolne odwzorowanie liniowe
$T_pM\to T_qN$ jest różniczką odpowiedniej lokalnej funkcji
w mapach. Dla $k\ge2$ dwa przedłużenia tego samego
$(k-1)$-dżetu różnią się tylko pochodną rzędu $k$.
W mapach oznaczmy różnicę najwyższych pochodnych
przez $B=D^kF-D^kG$. Po zmianie map $\varphi,\psi$
różnica przyjmuje postać
\[
 \widetilde B(v_1,\ldots,v_k)
 =D\psi\,B(D\varphi^{-1}v_1,\ldots,D\varphi^{-1}v_k),
\]
gdzie różniczki liczymy w odpowiednich punktach
bazowych. Jedyny wyraz z pochodną rzędu $k$
powstaje, gdy wszystkie $k$ różniczkowań padają
na $F$ z pierwszymi pochodnymi zmian map.
Pozostałe składniki zależą tylko od niższych
pochodnych, więc skracają się przy odejmowaniu.
Otrzymujemy regułę transformacji elementu
$S^kT_p^*M\otimes T_qN$. Dodanie dowolnego takiego
tensora do najwyższej pochodnej daje nowy dżet,
realizowany lokalnie przez wielomian Taylora.
Działanie jest wolne i przechodnie, co dowodzi
struktury afinicznej bez wyboru wyróżnionego zera.
\end{proof}

\begin{uwaga}
Nie należy nazywać całej $J^k(M,N)$ wiązką wektorową
nad $M\times N$ dla $k\ge2$: zmiany współrzędnych
zawierają wyrazy niższych rzędów. Różnica między
\emph{wiązka wektorowa} a \emph{wiązka afiniczna}
jest w tym miejscu istotna.
\end{uwaga}

\section{Dżety przekrojów, holonomiczność i formy kontaktowe}
Niech $\pi:E\to M$ będzie gładką wiązką lokalnie trywialną.
Dżety lokalnych przekrojów $s$ tworzą $J^kE$.
W lokalnej trywializacji współrzędne mają postać
$(x^i,u^a,u^a_\alpha)$ z $1\le|\alpha|\le k$.
Oznaczamy też $u^a_0=u^a$. Przekroje klasy $C^k$
mają dżety w tej samej przestrzeni, ponieważ ich
skończone dane pochodnych w ustalonym punkcie
można zrealizować przez gładki lokalny przekrój
zadany wielomianem w mapie włókna.
Każde $j^ks$ jest przekrojem $J^kE\to M$.

\begin{definicja}[Przekrój holonomiczny]
\label{def:dzety-holonomic}
Przekrój ciągły $\sigma:U\to J^kE$ nazywamy
\emph{holonomicznym}, jeżeli istnieje przekrój
$s:U\to E$ klasy $C^k$ taki, że $\sigma=j^ks$.
Gdy ponadto $\sigma$ jest klasy $C^1$, to $s$ jest
klasy $C^{k+1}$; wynika to kolejno ze współrzędnych
$u^a_\alpha\circ\sigma$, które są pochodnymi $s$.
\end{definicja}

Na $J^kE$ występują \emph{formy kontaktowe}
\[
 \theta^a_\alpha=\mathrm d u^a_\alpha
 -\sum_{i=1}^n u^a_{\alpha+e_i}\,\mathrm d x^i,
 \qquad |\alpha|\le k-1.
\]
Dla $\alpha=0$ przyjmujemy $u_0^a=u^a$.

\begin{twierdzenie}[Kryterium holonomiczności]
\label{thm:dzety-contact}
Niech $U\subseteq M$ będzie otwarty, $k\ge1$,
a $\sigma:U\to J^kE$ będzie przekrojem klasy $C^1$.
Wtedy $\sigma=j^ks$ dla pewnego przekroju
$s:U\to E$ klasy $C^{k+1}$ wtedy i tylko wtedy, gdy
$\sigma^*\theta^a_\alpha=0$ dla wszystkich form
kontaktowych, w dowolnych lokalnych trywializacjach.
\end{twierdzenie}
\begin{proof}
Dla $\sigma=j^ks$ zachodzi
$u^a_\alpha\circ\sigma=\partial^\alpha s^a$,
więc $\sigma^*\theta^a_\alpha=0$ wynika z różniczkowania.
Odwrotnie, niech $s=\pi_0^k\circ\sigma$, tj. składowa
zerowego rzędu. Równości $\sigma^*\theta^a_0=0$ dają
$\partial_i s^a=u^a_{e_i}\circ\sigma$. Indukcyjnie,
$\sigma^*\theta^a_\alpha=0$ implikuje
$\partial_i(u^a_\alpha\circ\sigma)
=u^a_{\alpha+e_i}\circ\sigma$. Otrzymujemy
$u^a_\alpha\circ\sigma=\partial^\alpha s^a$
dla $|\alpha|\le k$. Ponieważ wszystkie współrzędne $u^a_\alpha\circ\sigma$
są klasy $C^1$, otrzymujemy ponadto $s\in C^{k+1}$.
W szczególności regularność $C^1$ pola dżetów nie jest
tożsama z regularnością $C^1$ przekroju wyjściowego.
Konstrukcja jest zgodna na nakładaniu się map.
\end{proof}

\begin{przyklad}[Formalne dane nie muszą być pochodnymi]
W trywialnej wiązce $\mathbb R\times\mathbb R\to\mathbb R$
przekrój $\sigma(x)=(x,u(x),p(x))=(x,0,1)$
przestrzeni $J^1E$ ma zerową składową wartości i
stałą składową pochodnej równą $1$. Nie pochodzi od
żadnej funkcji $u$, ponieważ
$\sigma^*(\mathrm du-p\,\mathrm dx)=-\mathrm dx\ne0$.
\end{przyklad}

\begin{center}
\begin{tikzpicture}[>=stealth, every node/.style={font=\small}]
 \node (jk) at (0,1) {$J^kE$};
 \node (jm) at (2.2,1) {$J^{k-1}E$};
 \node (je) at (3.8,1) {$\cdots$};
 \node (e) at (5.2,1) {$E$};
 \node (m) at (6.6,1) {$M$};
 \node (u) at (0,-.25) {$U\subset M$};
 \draw[->] (jk) -- node[above] {$\pi_{k-1}^k$} (jm);
 \draw[->] (jm) -- (je);
 \draw[->] (je) -- (e);
 \draw[->] (e) -- node[above] {$\pi$} (m);
 \draw[->] (u) -- node[left] {$j^k s$} (jk);
\end{tikzpicture}

{\small Rzutowania w wieży dżetów odrzucają pochodne najwyższego rzędu.}
\end{center}

\section{Równania różniczkowe jako podzbiory przestrzeni dżetów}
\begin{definicja}[Równanie różniczkowe rzędu $k$]
\label{def:dzety-pde}
Niech $\pi:E\to M$ będzie gładką wiązką.
Układ równań różniczkowych rzędu nie większego niż $k$
możemy opisać \emph{lokalnie} zbiorem zer
$\mathcal R\cap V=\{z\in V:F(z)=0\}$, gdzie
$V\subseteq J^kE$ jest otwarty, a
$F:V\to\mathbb R^r$ jest gładkie. Nie każda globalna
podrozmaitość równań musi być zbiorem zer jednego
globalnie określonego odwzorowania do $\mathbb R^r$.
Przekrój $s$ klasy co najmniej $C^k$ jest rozwiązaniem,
gdy $j^ks(x)\in\mathcal R$ dla każdego $x$ z jego
dziedziny określenia.
\end{definicja}

\begin{przyklad}[Hamilton--Jacobi i równanie Laplace'a]
Dla rzeczywistej funkcji $u$ na $\mathbb R^n$
równanie Hamiltona--Jacobiego
$H(x,\nabla u(x))=0$ jest podzbiorem $J^1E$
opisanym przez $H(x,p)=0$. Równanie Laplace'a
$\sum_i\partial_{ii}u=0$ jest hiperpłaszczyzną
w każdym włóknie $J^2E\to J^1E$ o równaniu
$\sum_i u_{2e_i}=0$. W obu przypadkach sam wybór
punktu z $\mathcal R$ nie wystarcza: rozwiązanie
musi ponadto wyznaczać przekrój holonomiczny.
\end{przyklad}

\begin{definicja}[Symbol główny operatora liniowego]
\label{def:dzety-symbol}
Dla skalarnego liniowego operatora
$Lu=\sum_{|\alpha|\le k}a_\alpha(x)\partial^\alpha u$
o gładkich współczynnikach, na otwartym $U\subset\mathbb R^n$,
jego symbol główny jest wielomianem
$\sigma_k(L)(x,\xi)=\sum_{|\alpha|=k}a_\alpha(x)\xi^\alpha$.
\end{definicja}

\begin{uwaga}
Symbol główny jest zależny wyłącznie od najwyższej części
dżetu; pod zmianą map jego współczynniki przekształcają się
zgodnie z tensorową częścią struktury afinicznej
Twierdzenia~\ref{thm:dzety-affine}. Dla Laplasjanu
$\sigma_2(\Delta)(x,\xi)=|\xi|^2$.
\end{uwaga}

\begin{przyklad}[Równanie Eulera--Lagrange'a]
Niech $\Omega\subset\mathbb R^n$ będzie ograniczonym obszarem,
$u\in C^2(\Omega)$, a $L(x,z,p)$ będzie funkcją klasy $C^2$.
Zakładamy, że $L(x,u(x),\nabla u(x))$ jest całkowalna na $\Omega$.
Funkcjonał zależący od $1$-dżetu ma postać
\[
 \mathcal A(u)=\int_\Omega L(x,u(x),\nabla u(x))
 \,\mathrm d\lambda_n(x).
\]
Dla $\varphi\in C_c^\infty(\Omega)$ i zróżniczkowanego
pod całką wyrażenia (na zwartym nośniku $\varphi$) mamy
\[
 \frac{\mathrm d}{\mathrm dt}\mathcal A(u+t\varphi)\big|_{t=0}
 =\int_\Omega\left(L_z\varphi+
 \sum_i L_{p_i}\partial_i\varphi\right)\mathrm d\lambda_n(x).
\]
Jeżeli $u$ jest punktem stacjonarnym, czyli powyższa
pierwsza wariacja znika dla wszystkich
$\varphi\in C_c^\infty(\Omega)$, to po całkowaniu
przez części i zastosowaniu lematu podstawowego rachunku
wariacyjnego otrzymujemy równanie
$L_z-\sum_i\partial_i L_{p_i}=0$ w $\Omega$.
Jest ono zazwyczaj równaniem drugiego rzędu, a zatem
jego geometria znajduje się w $J^2E$, mimo że
sam funkcjonał zależy tylko od $J^1E$.
\end{przyklad}

\section{Od gładkości do miary: gęstość i granice przybliżone}
Odtąd pracujemy na $\mathbb R^n$. Przez $\lambda_n$ oznaczamy
miarę Lebesgue'a, a przez $B(x,r)$ kulę otwartą.

\begin{definicja}[Punkt gęstości i granica przybliżona]
\label{def:dzety-ap-limit}
Dla zbioru mierzalnego $A\subset\mathbb R^n$ punkt $x$
ma gęstość $1$ względem $A$, jeżeli
\[
 \lim_{r\downarrow0}
 \frac{\lambda_n(A\cap B(x,r))}{\lambda_n(B(x,r))}=1.
\]
Jeżeli $v$ jest funkcją mierzalną w otoczeniu $x$,
to piszemy $\operatorname{ap}\!\lim_{y\to x}v(y)=L$,
gdy dla każdego $\varepsilon>0$ zbiór
$\{y:|v(y)-L|>\varepsilon\}$ ma w $x$ gęstość $0$.
W przypadku funkcji określonej tylko na zbiorze mierzalnym
$A$ zakładamy, że $x$ jest punktem gęstości $1$ zbioru $A$
i liczymy gęstość zbiorów wyjątkowych w $A\cap B(x,r)$.
\end{definicja}

\begin{definicja}[Przybliżony $k$-dżet]
\label{def:dzety-approx-jet}
Niech $A\subset\mathbb R^n$ będzie mierzalny,
$x\in A$ punktem jego gęstości $1$, a
$f:A\to\mathbb R^m$ funkcją mierzalną.
Mówimy, że $f$ ma \emph{przybliżony dżet Peana rzędu $k\ge1$}
w $x$, jeżeli istnieje wielomian wektorowy
$P_x(h)=\sum_{|\alpha|\le k}a_\alpha(x)h^\alpha/\alpha!$,
przy $a_0(x)=f(x)$, taki że
\[
 \operatorname{ap}\!\lim_{\substack{y\to x\\y\in A}}
 \frac{f(y)-P_x(y-x)}{|y-x|^k}=0.
\]
Dla $k=1$ współczynniki liniowe wyznaczają
\emph{różniczkę przybliżoną} $D_{\mathrm{ap}}f(x)$.
W tym sformułowaniu wartość $f(x)$ jest częścią danych:
zmiana wartości funkcji w samym punkcie $x$ może więc
zniszczyć istnienie dżetu w tym punkcie. Natomiast
\emph{prawie wszędzie} pojęcie to zależy wyłącznie od
klasy funkcji modulo zbiory miary zero: dwa równe prawie
wszędzie przedstawiciele mają te same dżety w punktach
gęstości zbioru, na którym są równe.
\end{definicja}

\begin{lemat}[Jednoznaczność przybliżonego dżetu]
\label{lem:dzety-approx-unique}
Jeżeli $x$ jest punktem gęstości $1$ zbioru $A$,
to wielomian $P_x$ spełniający
Definicję~\ref{def:dzety-approx-jet} jest jednoznaczny.
\end{lemat}
\begin{proof}
Niech $Q=P_x-\widetilde P_x$ będzie różnicą dwóch takich
wielomianów. Wtedy $Q(0)=0$ i
$\operatorname{ap}\!\lim_{h\to0}Q(h)/|h|^k=0$.
Gdyby $Q\ne0$, jego pierwszy niezerowy jednorodny składnik
$Q_j$, $j\le k$, byłby niezerowy na pewnym otwartym
podzbiorze sfery jednostkowej. Na stożku o dodatnim kącie
bryłowym zachodziłoby
$|Q(h)|/|h|^k\ge c|h|^{j-k}$ dla dostatecznie małych $h$.
Na mniejszym otwartym fragmencie sfery $|Q_j|\ge2c>0$,
a wyższe składniki po podzieleniu przez $|h|^j$
są jednostajnie $O(|h|)$. Stąd dla dostatecznie
małych $|h|\le1$ iloraz jest nie mniejszy niż $c$.
Stożek $C$ ma względną miarę $\gamma>0$ w każdej
kuli wokół zera. Ponieważ $A-x$ ma gęstość $1$,
\[
 \liminf_{r\downarrow0}
 \frac{\lambda_n(C\cap(A-x)\cap B(0,r))}
 {\lambda_n(B(0,r))}\ge\gamma.
\]
Istotnie, od względnej miary stożka odejmujemy
co najwyżej względną miarę dopełnienia $A-x$,
która dąży do zera. Z nierówności trójkąta
zbiór, na którym $|Q(h)|/|h|^k>c/2$ w $A-x$,
zawiera się w sumie dwóch zbiorów wyjątkowych
dla obu wielomianów z progiem $c/4$.
Ich gęstości są zerowe, co daje sprzeczność.
\end{proof}

\begin{przyklad}[Funkcja nieciągła z przybliżoną pochodną]
\label{ex:dzety-approx}
Niech $A=\bigcup_{j\ge1}(2^{-j},2^{-j}+2^{-3j})$
i $f(x)=\chi_A(x)$, gdzie $\chi_A$ oznacza funkcję
charakterystyczną zbioru $A$ i $f(0)=0$. Wtedy $A$ ma
w zerze gęstość $0$, zatem $f$ ma w zerze przybliżoną
granicę równą $0$. Ponadto
$\lambda_1(A\cap(-r,r))=O(r^3)$ dla $r\downarrow0$,
a więc również $D_{\mathrm{ap}}f(0)=0$.
Zwykła ciągłość (a tym bardziej różniczkowalność) nie zachodzi,
ponieważ w każdym otoczeniu zera są punkty $f=1$.
\end{przyklad}

\begin{przyklad}[Dlaczego trzeba ustalić przedstawiciela?]
Niech $f=\chi_{\mathbb Q}$ na $\mathbb R$. W punktach
niewymiernych funkcja ma przybliżony dżet każdego rzędu
równy zeru, ponieważ zbiór $\mathbb Q$ ma miarę zero.
W punktach wymiernych $f(x)=1$, więc nie ma nawet
przybliżonej granicy równej własnej wartości $f(x)$.
Przedstawiciel zerowy tej samej klasy w $L^0$ ma natomiast
dżet zerowy w każdym punkcie. Rozróżnienie to nie ma
znaczenia dla twierdzeń prawie wszędzie.
\end{przyklad}

\begin{uwaga}[Gdzie tkwi istota mierzalności?]
Równość $D_{\mathrm{ap}}f(x)=A$ nie oznacza
zbieżności po \emph{wszystkich} punktach $y\to x$.
Wystarcza odrzucanie zbioru o gęstości zerowej,
który może być gęsty topologicznie. Dlatego dżety
przybliżone nie są dżetami kiełków w sensie
Definicji~\ref{def:dzety-klasa}.
\end{uwaga}

\section{Rademacher, Sobolew i $BV$}
\begin{twierdzenie}[Rademacher i mierzalny pierwszy dżet]
\label{thm:dzety-rademacher}
Niech $U\subset\mathbb R^n$ będzie otwarty, a
$f:U\to\mathbb R^m$ lokalnie Lipschitzowskie.
Wtedy $f$ jest różniczkowalne w $\lambda_n$-prawie każdym
punkcie $U$. W tych punktach istnieje przybliżony
$1$-dżet i $D_{\mathrm{ap}}f(x)=Df(x)$.
Ponadto $x\mapsto Df(x)$ ma mierzalny przedstawiciel.
\end{twierdzenie}
\begin{proof}
Istnienie zwykłej różniczki prawie wszędzie jest treścią
twierdzenia Rademachera. Zwykła granica implikuje
przybliżoną, więc równość dżetów jest natychmiastowa.
Zbiór różniczkowalności jest borelowski. Istotnie,
na zbiorze istnienia wszystkich pochodnych cząstkowych
ich wartości są granicami ilorazów różnicowych z
przyrostami $1/j$, a więc są funkcjami borelowskimi.
Warunek
$|f(x+h)-f(x)-A(x)h|\le |h|/m$
dla wszystkich dostatecznie małych niezerowych $h$
można, dzięki ciągłości $f$, sprawdzać na przeliczalnym
zbiorze przyrostów wymiernych. Kwantyfikatory
$\forall m\in\mathbb N\,\exists j\in\mathbb N\,
\forall h\in\mathbb Q^n\setminus\{0\}$ z $|h|<1/j$
i $x+h\in U$ opisują zatem borelowski podzbiór $S$,
który jest dokładnie zbiorem różniczkowalności.
Wynika to z ciągłości wyrażenia względem przyrostu $h$
poza $h=0$. Przedłużenie macierzy
$Df$ przez zero na dopełnienie tego zbioru jest mierzalne.
\end{proof}

\begin{przyklad}[Punkt o dwóch geometrycznych dżetach]
Funkcja $f(x)=|x|$ na $\mathbb R$ ma zwykłą różniczkę
$Df(x)=\operatorname{sgn}(x)$ dla $x\ne0$.
W zerze nie ma ani zwykłego, ani przybliżonego
$1$-dżetu: dodatnia i ujemna półprosta mają dodatnią
gęstość i dawałyby sprzeczne pochodne. Brak dżetu w jednym
punkcie nie narusza twierdzenia Rademachera.
\end{przyklad}

\begin{twierdzenie}[Dżet przybliżony funkcji Sobolewa]
\label{thm:dzety-sobolev}
Niech $U\subset\mathbb R^n$ będzie otwarty,
$1\le p\le\infty$, a $u\in W_{\mathrm{loc}}^{1,p}(U)$.
Wówczas istnieje przedstawiciel $u$, który ma
przybliżony $1$-dżet w $\lambda_n$-prawie każdym
punkcie $U$, przy czym
$D_{\mathrm{ap}}u(x)=\nabla u(x)$ prawie wszędzie.
Po prawej stronie występuje gradient słaby.
\end{twierdzenie}
\begin{proof}
Korzystamy z następującego twierdzenia o precyzyjnych
przedstawicielach Sobolewa: funkcja klasy $W^{1,1}_{\rm loc}$
jest przybliżenie różniczkowalna prawie wszędzie, a jej
różniczka przybliżona pokrywa się z gradientem słabym.
Standardowy dowód tego wyniku opiera się na różniczkowaniu
całek i przybliżeniu funkcji Sobolewa przez funkcje gładkie.
Dla $p>1$ można go uzyskać także z oszacowań funkcji
maksymalnej Hardy'ego--Littlewooda; przypadek $p=1$
wymaga wersji słabego typu $(1,1)$. Ponieważ
$W^{1,p}_{\rm loc}\subset W^{1,1}_{\rm loc}$,
otrzymujemy tezę. Podany rezultat o przedstawicielach
Sobolewa jest w tym dowodzie \emph{jawnie użytym
wcześniejszym twierdzeniem}, a nie wynikiem wyprowadzonym
z samej definicji dżetu.
\end{proof}

\begin{uwaga}[Funkcje $BV$ i punkty skoku]
\label{rem:dzety-bv}
Dla $u\in BV_{\rm loc}(U)$ istnieje przybliżona
różniczka w prawie każdym punkcie względem $\lambda_n$,
lecz pochodna dystrybucyjna jest miarą wektorową, która
może mieć część osobliwą. W pełnym twierdzeniu
strukturalnym dla $BV$ rozkłada się ją na
\[
 Du=\nabla u\,\lambda_n
 +(u^+-u^-)\nu_u\,\mathcal H^{n-1}\llcorner J_u+D^cu,
\]
gdzie $J_u$ jest przybliżonym zbiorem skoków,
$u^\pm$ jednostronnymi śladami, $\nu_u$ normalną,
a $D^cu$ częścią Cantora. Jest to odrębny, głębszy wynik:
nie należy wyprowadzać go wyłącznie z istnienia
przybliżonego $1$-dżetu prawie wszędzie.
\end{uwaga}

\begin{przyklad}[Skok a przybliżony dżet]
Dla $E=\{x\in\mathbb R^n:x_1>0\}$ funkcja
$u=\chi_E$ jest lokalnie stała poza hiperpłaszczyzną
$x_1=0$, więc $D_{\mathrm{ap}}u=0$ prawie wszędzie.
Jednak $Du=e_1\,\mathcal H^{n-1}\llcorner\{x_1=0\}$
(z orientacją odpowiadającą przejściu z $0$ do $1$).
Zatem przybliżony dżet pierwszego rzędu nie opisuje
całej miary pochodnej funkcji $BV$.
\end{przyklad}

\section{Jakobiany i wzory area/coarea odczytywane z dżetów}
Niech $f:U\subseteq\mathbb R^k\to\mathbb R^N$,
$k\le N$, będzie lokalnie Lipschitzowskie.
W punkcie różniczkowalności określamy
$J_kf(x)=\sqrt{\det(Df(x)^{\mathsf T}Df(x))}$.
Jest to funkcja samego \emph{liniowego składnika $1$-dżetu}.
Dla mierzalnego $A\subset U$ i nieujemnej
mierzalnej funkcji $h:U\to[0,\infty]$ wzór area ma postać
\[
 \int_A h(x)J_kf(x)\,\mathrm d\lambda_k(x)
 =\int_{\mathbb R^N}\sum_{x\in A\cap f^{-1}(\{y\})}
 h(x)\,\mathrm d\mathcal H^k(y).
\]
Dla lokalnie Lipschitzowskiego $F:V\subset\mathbb R^n\to
\mathbb R^m$, $n\ge m$, gdzie $V$ jest otwarty,
i nieujemnej mierzalnej $g:V\to[0,\infty]$, wzór
koobszaru brzmi
\[
 \int_V g(x)\sqrt{\det(DF(x)DF(x)^{\mathsf T})}
 \,\mathrm d\lambda_n(x)
 =\int_{\mathbb R^m}\left(
 \int_{F^{-1}(\{y\})}g(x)\,\mathrm d\mathcal H^{n-m}(x)
 \right)\mathrm d\lambda_m(y).
\]
W obu wzorach miary Hausdorffa są znormalizowane tak,
że $\mathcal H^j|_{\mathbb R^j}=\lambda_j$ dla całkowitych
$j\ge1$ oraz $\mathcal H^0$ jest miarą liczącą.
Dla zbiorów mierzalnych w sensie Lebesgue'a konieczne
są właściwe sformułowania twierdzeń, uwzględniające
ukończenie miar i mierzalność funkcji krotnościowej.
W szczególności sama teoria dżetów nie rozwiązuje
kwestii mierzalności włókien i całek po nich.

\begin{uwaga}
Teoria dżetów wyjaśnia, \emph{jakie dane} wyznaczają
lokalny współczynnik geometryczny. Sam wzór area/coarea
wymaga jednak niezależnie twierdzeń o mierzalności,
różniczkowaniu i krotnościach włókien. Nie jest więc
czysto formalną konsekwencją teorii dżetów.
\end{uwaga}

\section{Pola dżetów i twierdzenie Whitneya}
\begin{definicja}[Pole dżetów Whitneya klasy $C^k$]
\label{def:dzety-whitney-field}
Niech $K\subset\mathbb R^n$ będzie domknięty,
$k\in\mathbb N_0$, a dla każdego $|\alpha|\le k$
niech $a_\alpha:K\to\mathbb R$ będzie funkcją.
Kładziemy $P_x(z)=\sum_{|\alpha|\le k}
 a_\alpha(x)(z-x)^\alpha/\alpha!$.
Rodzinę $(P_x)_{x\in K}$ nazywamy
\emph{polem dżetów Whitneya klasy $C^k$}, jeżeli
wszystkie $a_\alpha$ są ciągłe i dla każdego zwartego
$K_0\subset K$ oraz każdego $|\alpha|\le k$ zachodzi
\[
 \lim_{\delta\downarrow0}
 \sup_{\substack{x,y\in K_0,\;0<|x-y|<\delta}}
 \frac{|D^\alpha P_x(y)-a_\alpha(y)|}
 {|x-y|^{k-|\alpha|}}=0,
\]
przy umowie, że supremum po zbiorze pustym wynosi $0$.
\end{definicja}

\begin{uwaga}
Dla $|\alpha|=k$ mianownik jest równy $1$:
warunek mówi o ciągłości najwyższych pochodnych.
Dla $k=1$ warunki na zwartym $K$ sprowadzają się do
ciągłości $a_0,v=(a_{e_1},\ldots,a_{e_n})$ i
$|a_0(y)-a_0(x)-v(x)\cdot(y-x)|=o(|x-y|)$
jednostajnie dla $x,y\in K$.
\end{uwaga}

\begin{lemat}[Kostki Whitneya i rozkład jedności]
\label{lem:dzety-whitney-cubes}
Niech $K\subset\mathbb R^n$ będzie niepustym zbiorem
domkniętym, różnym od całej przestrzeni, a
$\Omega=\mathbb R^n\setminus K$. Istnieje przeliczalna
rodzina domkniętych kostek diadycznych $Q\subset\Omega$,
o parami rozłącznych wnętrzach, których suma jest
równa $\Omega$, taka że
$\ell(Q)\asymp\operatorname{dist}(Q,K)$. Istnieje
również gładki rozkład jedności $(\varphi_Q)$ na $\Omega$
o nośnikach w nieznacznie powiększonych kostkach $Q^*$,
takich że
$|D^\alpha\varphi_Q|\le C_\alpha\ell(Q)^{-|\alpha|}$.
W każdym punkcie co najwyżej $N(n)$ funkcji
$\varphi_Q$ jest niezerowych.
\end{lemat}
\begin{proof}
\emph{Wybór kostek.}
Przez kostkę diadyczną rozumiemy domknięcie zbioru
$2^{-j}(m+[0,1)^n)$, gdzie $j\in\mathbb Z$,
$m\in\mathbb Z^n$. Wybieramy maksymalne względem
inkluzji kostki spełniające
\[
 8\sqrt n\,\ell(Q)\le\operatorname{dist}(Q,K).
\]
Dla $z\in\Omega$ niech $d=\operatorname{dist}(z,K)>0$.
Każda dostatecznie mała kostka zawierająca $z$
spełnia warunek wyboru, ponieważ jej odległość od $K$
jest nie mniejsza niż $d-\sqrt n\,\ell(Q)$.
Wśród kolejnych kostek rodzicielskich warunek
w końcu przestaje zachodzić: każda zawiera $z$,
więc jej odległość od $K$ jest nie większa niż $d$,
a długości boków rosną nieograniczenie.
Ostatnia dopuszczalna kostka jest maksymalna.
Dwie kostki diadyczne o przecinających się wnętrzach
są zagnieżdżone, więc różne kostki maksymalne
mają rozłączne wnętrza. Każda leży w $\Omega$,
a wszystkie razem pokrywają $\Omega$.

Jeśli $Q'$ jest kostką rodzicielską dla wybranej $Q$,
to $\ell(Q')=2\ell(Q)$, a z maksymalności
$\operatorname{dist}(Q',K)<16\sqrt n\,\ell(Q)$.
Każdy punkt $Q'$ jest odległy od $Q$ co najwyżej
o średnicę $Q'$. Stąd
\[
 8\sqrt n\,\ell(Q)\le\operatorname{dist}(Q,K)
 <18\sqrt n\,\ell(Q).
\]

\emph{Powiększenia i krotność pokrycia.}
Niech $Q^*$ będzie otwartą kostką o tym samym środku,
o boku $(5/4)\ell(Q)$. Dla $z\in\overline{Q^*}$
odległość od $Q$ jest co najwyżej
$\sqrt n\,\ell(Q)/8$. Z dolnego oszacowania
odległości $Q$ od $K$ oraz nierówności trójkąta mamy
$\operatorname{dist}(z,K)\ge(8-1/8)\sqrt n\,\ell(Q)$.
Z drugiej strony, dla dowolnego $q\in Q$
mamy $|z-q|\le(9/8)\sqrt n\,\ell(Q)$.
Biorąc kres dolny odległości $|q-b|$, $b\in K$,
uzyskujemy wygodne oszacowania
\[
 7\sqrt n\,\ell(Q)\le\operatorname{dist}(z,K)
 \le20\sqrt n\,\ell(Q)\qquad(z\in\overline{Q^*}).
\]
W szczególności $\overline{Q^*}\subset\Omega$.
Jeżeli powiększenia dwóch kostek zawierają ten
sam punkt, stosunek ich boków mieści się między
$7/20$ a $20/7$.

Ustalmy $z\in\Omega$ i $d=\operatorname{dist}(z,K)$.
Wszystkie kostki, których powiększenia zawierają $z$,
mają boki co najmniej $d/(20\sqrt n)$ i co najwyżej
$d/(7\sqrt n)$. Każda z nich leży w $B(z,d)$,
gdyż jej punkty są odległe od $z$ najwyżej
o $(9/8)\sqrt n\,\ell(Q)<d$.
Z rozłączności wnętrz, dla dowolnej skończonej
podrodziny liczącej $N$ kostek otrzymujemy
\[
 N\left(\frac{d}{20\sqrt n}\right)^n
 \le\lambda_n(B(z,d))=\lambda_n(B(0,1))d^n.
\]
Daje to ograniczenie liczności całej rodziny
przez stałą zależną tylko od $n$.
Ponadto pokrycie powiększeniami jest lokalnie
skończone: jeżeli zwarty $L\subset\Omega$ jest
niepusty, to odległość jego punktów od $K$
ma dodatnie minimum i skończone maksimum.
Każde $\overline{Q^*}$ przecinające $L$ ma więc
bok ograniczony z dołu i z góry, a środek leżący
w pewnym ograniczonym zbiorze. Rozłączne wnętrza
i dodatnie dolne ograniczenie objętości kostek
dają skończenie wiele takich kostek.

\emph{Rozkład jedności.}
Wybieramy ustaloną gładką funkcję obcięcia
$0\le\eta\le1$, równą $1$ na kostce jednostkowej
o środku zero, o zwartym nośniku w jej
powiększeniu o czynnik $5/4$. Dla środka $c_Q$
definiujemy $\eta_Q(z)=\eta((z-c_Q)/\ell(Q))$.
Wtedy
$|D^\alpha\eta_Q|\le C_\alpha\ell(Q)^{-|\alpha|}$.
Suma $S=\sum_Q\eta_Q$ jest lokalnie skończona
i gładka, a pokrycie przez $Q$ oraz ograniczona
krotność dają $1\le S\le N(n)$.
Kładziemy $\varphi_Q=\eta_Q/S$.
Są to nieujemne funkcje gładkie o sumie $1$,
z nośnikami zawartymi w $Q^*$.

Na $Q^*$ pochodne $S$ są ograniczone przez
$C_\alpha'\ell(Q)^{-|\alpha|}$, ponieważ w każdym
punkcie jedynie ograniczona liczba składników
może mieć tam niezerowe pochodne, a ich boki
są porównywalne z $\ell(Q)$ (stosujemy oszacowania
również na domknięciach powiększeń).
Z tożsamości $S\cdot S^{-1}=1$ i reguły Leibniza,
dla $|\alpha|>0$,
\[
 D^\alpha(S^{-1})
 =-S^{-1}\sum_{0<\beta\le\alpha}
 \binom{\alpha}{\beta}D^\beta S\,
 D^{\alpha-\beta}(S^{-1}).
\]
Indukcja względem $|\alpha|$, zaczynając od
$|S^{-1}|\le1$, daje na $Q^*$
$|D^\alpha(S^{-1})|\le C_\alpha''\ell(Q)^{-|\alpha|}$.
Reguła Leibniza dla $\eta_Q S^{-1}$ daje żądane
oszacowanie pochodnych $\varphi_Q$; poza $Q^*$
funkcja i wszystkie jej pochodne znikają.
\end{proof}

\begin{lemat}[Sklejanie ciągłych pochodnych przez zbiór domknięty]
\label{lem:dzety-glue}
Niech $K\subset\mathbb R^n$ będzie domknięty. Załóżmy, że
$H:\mathbb R^n\to\mathbb R$ jest ciągła, klasy $C^1$ poza $K$,
a $V:\mathbb R^n\to\mathbb R^n$ jest ciągła i spełnia
$V=\nabla H$ poza $K$. Jeżeli dla każdego $a\in K$
zachodzi $H(a+h)=H(a)+V(a)\cdot h+o(|h|)$, to
$H\in C^1(\mathbb R^n)$ i $\nabla H=V$ wszędzie.
\end{lemat}
\begin{proof}
Poza $K$ teza jest znana. W punkcie $a\in K$ podany
warunek jest definicją różniczkowalności i daje
$\nabla H(a)=V(a)$. Ponieważ $V$ jest ciągła na całej
przestrzeni, otrzymany gradient jest ciągły, czyli
$H\in C^1$. Nie trzeba różniczkować po krzywych
zawartych w $K$.
\end{proof}

\begin{twierdzenie}[Twierdzenie Whitneya o przedłużaniu]
\label{thm:dzety-whitney-extension}
Niech $K\subset\mathbb R^n$ będzie domknięty,
$k\in\mathbb N_0$, a $(a_\alpha)_{|\alpha|\le k}$
polem dżetów Whitneya klasy $C^k$ na $K$.
Wtedy istnieje $F\in C^k(\mathbb R^n)$ spełniająca
$D^\alpha F(x)=a_\alpha(x)$ dla wszystkich
$x\in K$ i $|\alpha|\le k$.
Warunki Definicji~\ref{def:dzety-whitney-field} są
zarazem konieczne dla takich przedłużeń.
\end{twierdzenie}

\begin{proof}
Konieczność: dla $F\in C^k(\mathbb R^n)$ kładziemy
$a_\alpha=D^\alpha F|_K$. Ustalmy zwarty $K_0\subset K$
i $|\alpha|\le k$, niech $m=k-|\alpha|$.
Dla $m=0$ żądane oszacowanie wynika z jednostajnej
ciągłości $D^\alpha F$ na $K_0$.
Dla $m\ge1$ stosujemy wzór Taylora z resztą całkową
do funkcji $t\mapsto D^\alpha F(x+t(y-x))$,
odejmując od reszty wyraz stopnia $m$ w punkcie $x$.
Dostajemy
\[
 a_\alpha(y)-D^\alpha P_x(y)
 =\sum_{|\gamma|=m}\frac{m(y-x)^\gamma}{\gamma!}
 \int_0^1(1-t)^{m-1}
 [D^{\alpha+\gamma}F(x+t(y-x))-D^{\alpha+\gamma}F(x)]
 \,\mathrm d\lambda_1(t).
\]
Dla $|x-y|\le1$ wszystkie punkty odcinka należą
do zwartego zbioru $\{z:\operatorname{dist}(z,K_0)\le1\}$.
Moduł ciągłości pochodnych rzędu $k$ na tym zbiorze
ogranicza więc iloraz reszty przez $|y-x|^m$
jednostajnie przez wielokrotność funkcji dążącej do zera.
Przedstawimy konstrukcję wystarczalności; dla $K=\varnothing$
można przyjąć $F=0$.

\emph{Krok 1: kostki Whitneya.}
Jeżeli $K=\mathbb R^n$, przyjmujemy $F=a_0$.
Dla $|\beta|<k$ warunki Whitneya w otoczeniu
dowolnego $x$ dają
$a_\beta(x+h)=a_\beta(x)+\sum_i a_{\beta+e_i}(x)h_i+o(|h|)$:
wyrazy stopni co najmniej $2$ są $O(|h|^2)$,
a reszta jest $o(|h|^{k-|\beta|})$.
Zatem $a_\beta$ jest różniczkowalna i jej pochodne
pierwszego rzędu są równe odpowiednim ciągłym
$a_{\beta+e_i}$. Indukcyjnie $a_0\in C^k$
i $D^\beta a_0=a_\beta$. Dla $k=0$ wystarcza
ciągłość $a_0$.
W przeciwnym przypadku Lemat~\ref{lem:dzety-whitney-cubes}
dostarcza rozkładu zbioru
$\Omega=\mathbb R^n\setminus K$
na kostki $Q$, wraz z rozkładem jedności
$(\varphi_Q)$ i oszacowaniami ich pochodnych.

\emph{Krok 2: sklejanie wielomianów.}
Dla każdej kostki wybierzmy $b_Q\in K$ takie, że
$\operatorname{dist}(b_Q,Q)=\operatorname{dist}(Q,K)$, i połóżmy
$F(z)=\sum_Q\varphi_Q(z)P_{b_Q}(z)$ dla $z\in\Omega$,
a na $K$ niech $F(x)=a_0(x)$. Suma jest lokalnie skończona
i funkcja jest gładka poza $K$.

\emph{Krok 3: oszacowanie w pobliżu brzegu.}
Ustalmy $x\in K$. Wszystkie punkty $a,b\in K$
potrzebne w oszacowaniach dla $z$ dostatecznie
bliskiego $x$ leżą w zwartym $K\cap\overline{B(x,1)}$,
jak wynika z poniższych oszacowań odległości.
Na tym zbiorze niech $\omega(t)$ będzie maksimum
po $|\beta|\le k$ supremów ilorazów z definicji
Whitneya dla $0<|a-b|\le t$, z wartością zero
dla pustej rodziny. Dla dostatecznie małych $t>0$
jest skończona, niemalejąca i dąży do zera
przy $t\downarrow0$. Dla $a,b$ z tego zbioru,
$0<|a-b|\le t$, mamy
\[
 |D^\beta P_a(b)-a_\beta(b)|
 \le\omega(t)|a-b|^{k-|\beta|},\qquad |\beta|\le k.
\]
Dla $a=b$ różnica wielomianów jest zerowa.
Dla $a\ne b$ rozwijamy wielomian $D^\beta(P_a-P_b)$
dokładnie w punkcie $b$:
\[
 D^\beta(P_a-P_b)(z)
 =\sum_{|\gamma|\le k-|\beta|}
 \frac{D^{\beta+\gamma}P_a(b)-a_{\beta+\gamma}(b)}
 {\gamma!}(z-b)^\gamma.
\]
Jeżeli $|z-a|+|z-b|\le Ct$, to $|a-b|\le Ct$.
Każdy składnik sumy jest ograniczony w module
przez stałą razy
$\omega(Ct)t^{k-|\beta|-|\gamma|}t^{|\gamma|}$.
Skończenie wiele składników daje
$|D^\beta(P_a-P_b)(z)|\le C'\omega(Ct)t^{k-|\beta|}$.
Niech teraz $z\in\Omega$, $r=\operatorname{dist}(z,K)$,
a $a\in K$ jest punktem najbliższym $z$, więc
$|z-a|=r$ (odległość od niepustego zbioru domkniętego
w $\mathbb R^n$ jest osiągana). Dla każdej kostki
$Q$ z $z\in\overline{Q^*}$ mamy $\ell(Q)\asymp r$,
a z wyboru $b_Q$ i oszacowań kostek
$|b_Q-z|\le C\ell(Q)\le C'r$. W konsekwencji
$|b_Q-a|\le C''r$. Ponieważ $r\le|z-x|$,
punkty $a,b_Q$ należą do ustalonego zwartego
fragmentu $K$ dla dostatecznie małego $|z-x|$.

Traktując $a$ jako ustalony punkt podczas
różniczkowania, z $\sum_Q\varphi_Q=1$ mamy
$F-P_a=\sum_Q\varphi_Q(P_{b_Q}-P_a)$ na $\Omega$.
Reguła Leibniza daje w punkcie $z$
\[
 D^\beta(F-P_a)(z)
 =\sum_{Q:\ z\in\overline{Q^*}}\sum_{\gamma\le\beta}
 \binom{\beta}{\gamma}
 D^\gamma\varphi_Q(z)\,
 D^{\beta-\gamma}(P_{b_Q}-P_a)(z).
\]
Każdy iloczyn jest ograniczony w module przez
$C r^{-|\gamma|}\omega(Cr)r^{k-|\beta|+|\gamma|}$.
Liczba kostek jest ograniczona przez stałą
wymiarową, a liczba wielowskaźników przez stałą
zależną od $n,k$. Ostatecznie dla $|\beta|\le k$
\[
 |D^\beta F(z)-D^\beta P_a(z)|
 \le C\omega(Cr)r^{k-|\beta|}.
\]
W szczególności pochodne do rzędu $k$ mają przy $K$
granice zadane przez współczynniki $a_\beta$.

\emph{Krok 4: regularność na $K$.}
Dla każdego $|\beta|\le k$ definiujemy na całej
przestrzeni funkcję $H_\beta$ równą $D^\beta F$
na $\Omega$ i $a_\beta$ na $K$. W szczególności
$H_0=F$. Ustalmy $x\in K$ i oznaczmy $d=|z-x|>0$.
Dla $z\in\Omega$ wybieramy $a$ jak wyżej;
$r\le d$ i $|a-x|\le2d$. Wobec tego
\[
 \begin{split}
 |H_\beta(z)-D^\beta P_x(z)|
 &\le |D^\beta F(z)-D^\beta P_a(z)|\\
 &\quad+|D^\beta(P_a-P_x)(z)|\\
 &\le C\omega(Cr)r^{k-|\beta|}
       +C'\omega(C'd)d^{k-|\beta|}\\
 &\le C''\omega(C''d)d^{k-|\beta|}.
 \end{split}
\]
Ostatnia nierówność wykorzystuje niemalejący
charakter $\omega$ i nieujemność $k-|\beta|$.
Dla $z\in K$ to samo oszacowanie wynika
bezpośrednio z warunku Whitneya dla $x,z$.
Zatem dla dowolnego kierunku zbliżania $z\to x$
mamy $H_\beta(z)-D^\beta P_x(z)=o(d^{k-|\beta|})$.
Dla $|\beta|=k$ interpretujemy $o(d^0)$ jako $o(1)$.
Ponieważ $D^\beta P_x(x)=a_\beta(x)$,
wszystkie $H_\beta$ są ciągłe w $x$; poza $K$
są gładkie. Dla $k=0$ kończy to dowód.

Jeżeli $|\beta|<k$, rozwijamy wielomian w punkcie $x$:
\[
 D^\beta P_x(x+h)
 =a_\beta(x)+\sum_i a_{\beta+e_i}(x)h_i+O(|h|^2).
\]
W połączeniu z poprzednim oszacowaniem, ponieważ
$k-|\beta|\ge1$, otrzymujemy
$H_\beta(x+h)=H_\beta(x)+\sum_i H_{\beta+e_i}(x)h_i+o(|h|)$.
Poza $K$ mamy już $\partial_i H_\beta=H_{\beta+e_i}$.
Wektor $(H_{\beta+e_i})_i$ jest ciągły na całej przestrzeni. Stosując
Lemat~\ref{lem:dzety-glue} kolejno do pochodnych
rzędów $0,1,\ldots,k-1$, wykazujemy
$F\in C^k(\mathbb R^n)$ oraz $D^\alpha F|_K=a_\alpha$.
W ten sposób przejście przez $K$ nie wymaga
nieuzasadnionego różniczkowania funkcji na samym zbiorze
$K$, który może mieć puste wnętrze.
\end{proof}

\begin{przyklad}[Dlaczego sama ciągłość dżetów nie wystarcza?]
Niech $K=\{0\}\cup\{1/j:j\ge1\}\subset\mathbb R$,
$a_0(0)=0$, $a_0(1/j)=1/j$ i $a_1\equiv0$.
Obie funkcje $a_0,a_1$ są ciągłe na $K$, ale
$|a_0(1/j)-a_0(0)-a_1(0)/j|/(1/j)=1$.
Nie istnieje więc funkcja $C^1$ o tych wartościach i
pochodnych. Jest to konkretne naruszenie warunku Whitneya.
\end{przyklad}

\section{Twierdzenie Whitneya--Luzina: od dżetów przybliżonych do $C^1$}
\begin{lemat}[Mierzalność różniczki przybliżonej]
\label{lem:dzety-ap-meas}
Niech $U\subset\mathbb R^n$ będzie otwarty,
a $f:U\to\mathbb R^m$ mierzalna. Na mierzalnym
zbiorze punktów, w których $f$ jest przybliżenie
różniczkowalna, odwzorowanie $D_{\mathrm{ap}}f$
ma mierzalny przedstawiciel.
\end{lemat}
\begin{proof}
Zmienimy wartości na zbiorze miary zero tylko wtedy,
gdy jest to potrzebne do zdefiniowania funkcji na całym $U$.
Dla $M\in\mathbb N$ i $r>0$ rozważmy macierz
\[
 T_{M,r}(x)=\frac{n}{\lambda_n(B(0,r))}
 \int_{B(0,r)}\frac{f(x+h)-f(x)}{|h|}
 \frac{h^{\mathsf T}}{|h|}
 \chi_{\{|f(x+h)-f(x)|\le M|h|\}}
 \,\mathrm d\lambda_n(h),
\]
gdzie dla $x+h\notin U$ całkowaną funkcję przyjmujemy za zero.
Jest to odwzorowanie mierzalne, ponieważ całkowana funkcja
jest mierzalna i ograniczona przez $M$.
Jeżeli $f$ ma różniczkę przybliżoną $A$ w $x$ i
$M>\|A\|+1$, to na kulach malejących do zera
iloraz różnicowy jest w mierze bliski
$A(h/|h|)$; obcięcie nie usuwa prawie wszystkich
kierunków. Ponieważ całkowana funkcja jest ograniczona,
własność zbieżności w mierze daje
$\lim_{r\downarrow0}T_{M,r}(x)=A$.
Wykorzystaliśmy tożsamość
$\lambda_n(B(0,r))^{-1}\int_{B(0,r)}
 hh^{\mathsf T}/|h|^2\,\mathrm d\lambda_n(h)=I_n/n$.
Dla ciągu $r=1/j$ granice macierzy $T_{M,1/j}$ są
mierzalne na zbiorach swojego istnienia. W każdym punkcie
przybliżonej różniczkowalności wszystkie dostatecznie
duże wartości $M$ dają tę samą granicę $A$.
Wybór pierwszego stabilnego ogona tych granic jest
operacją na przeliczalnej rodzinie funkcji mierzalnych,
a więc daje mierzalny przedstawiciel $A$.
\end{proof}

\begin{twierdzenie}[Whitney--Luzin dla pierwszego dżetu]
\label{thm:dzety-whitney-lusin}
Niech $U\subset\mathbb R^n$ będzie otwarty,
a $f:U\to\mathbb R^m$ funkcją mierzalną,
przybliżenie różniczkowalną w $\lambda_n$-prawie każdym
punkcie $U$. Dla każdego zbioru mierzalnego $A\subset U$
o skończonej mierze i każdego $\varepsilon>0$ istnieją
zbiór zwarty $K\subset A$ i odwzorowanie
$F\in C^1(\mathbb R^n,\mathbb R^m)$ takie, że
\[
 \lambda_n(A\setminus K)<\varepsilon,
 \qquad F(x)=f(x),\quad DF(x)=D_{\mathrm{ap}}f(x)
 \quad (x\in K).
\]
Odwrotnie, jeżeli dla każdego mierzalnego $A\subset U$
o skończonej mierze i każdego $\varepsilon>0$ istnieją
$K,F$ spełniające powyższą równość wartości (bez wymogu
dotyczącego pochodnych), to $f$ jest przybliżenie
różniczkowalna prawie wszędzie na $U$.
\end{twierdzenie}

\begin{proof}
Wystarczy dowodzić dla funkcji rzeczywistej: dla wektorowych
stosujemy ten sam wybór zwartego zbioru i wspólne
oszacowania normy w $\mathbb R^m$.

\emph{Krok 1: duży zbiór z ciągłymi danymi.}
Niech $D$ będzie pełnomiarowym podzbiorem $U$, gdzie
istnieje różniczka przybliżona $G(x)=D_{\mathrm{ap}}f(x)$.
Dla sformułowania „prawie wszędzie” wybieramy mierzalny
zbiór $D$ o pełnej mierze, na którym dżet istnieje;
na nim działa Lemat~\ref{lem:dzety-ap-meas}.
Funkcje $f,G$ są więc mierzalne na $D$. Z twierdzenia Luzina i regularności
miary Lebesgue'a wybieramy zwarty $K_0\subset A\cap D$,
$\lambda_n(A\setminus K_0)<\varepsilon/3$, tak że
$f|_{K_0}$ i $G|_{K_0}$ są ciągłe.
W szczególności $G|_{K_0}$ jest jednostajnie ciągła.
Ponieważ $K_0\Subset U$, ma dodatnią odległość od
$\mathbb R^n\setminus U$.

\emph{Krok 2: jednostajna kontrola zbiorów złych.}
Dla $j\in\mathbb N$ i $r>0$ oznaczmy przez
$\vartheta_{j,r}(x)$ iloraz
\[
 \frac{\lambda_n\bigl(\{z\in B(x,r)\cap U:
 |f(z)-f(x)-G(x)(z-x)|>j^{-1}|z-x|\}\bigr)}
 {\lambda_n(B(x,r))}.
\]
Z przybliżonej różniczkowalności wynika
$\vartheta_{j,r}(x)\to0$ przy $r\downarrow0$ dla
każdego $x\in D$. Dla ustalonego $j$ funkcje
$H_{j,l}(x)=\sup_{0<r<1/l}\vartheta_{j,r}(x)$ są
mierzalne: dla ustalonego $x$ licznik zależy ciągle
od $r$ (sfery mają miarę zero), toteż supremum można
brać po promieniach wymiernych; mierzalność całki
zależącej od $x$ daje twierdzenie Tonellego.
Mamy $H_{j,l}(x)\downarrow0$.

Wybierzmy stałą $\delta=8^{-n}$. Z twierdzenia Jegorowa
stosowanego kolejno dla $j=1,2,\ldots$ wybieramy
mierzalne podzbiory $E_j\subset K_0$ tak, że
$\lambda_n(K_0\setminus E_j)<\varepsilon/(3\cdot 2^j)$
oraz $H_{j,l}\to0$ jednostajnie na $E_j$.
Niech $E=\bigcap_j E_j$. Następnie wybierzmy zwarty
$K\subset E$ z $\lambda_n(E\setminus K)<\varepsilon/3$.
Wtedy $\lambda_n(A\setminus K)<\varepsilon$.
Dla każdego $j$ istnieje $\rho_j>0$ taka, że
$\vartheta_{j,r}(x)<\delta$ jednocześnie dla
$x\in K$ i $0<r<\rho_j$.

\emph{Krok 3: z kontroli w mierze do warunku Whitneya.}
Niech $x,y\in K$, $d=|x-y|>0$ i
$2d<\rho_j$; wymagamy także, aby $B(x,2d),B(y,2d)$
były zawarte w $U$. Przecięcie tych kul zawiera
$B((x+y)/2,3d/2)$, więc jego miara jest większa
niż suma dwóch górnych oszacowań
$\delta\lambda_n(B(0,2d))$ na zbiory złe
(dla $\delta=8^{-n}$). Istnieje zatem punkt $z$
w przecięciu, dla którego jednocześnie
\[
 |f(z)-f(x)-G(x)(z-x)|\le 2d/j,
 \quad |f(z)-f(y)-G(y)(z-y)|\le 2d/j.
\]
Odejmując równości otrzymujemy
\[
 |f(y)-f(x)-G(x)(y-x)|
 \le4d/j+2d\,\|G(x)-G(y)\|.
\]
Jednostajna ciągłość $G$ na $K$ i dowolność $j$
dowodzą, że iloraz lewej strony przez $|x-y|$
dąży jednostajnie do zera. Tak więc $(f|_K,G|_K)$
jest polem dżetów Whitneya klasy $C^1$.

\emph{Krok 4: przedłużenie.}
Twierdzenie~\ref{thm:dzety-whitney-extension}
daje $F\in C^1(\mathbb R^n)$ o danych wartościach
i różniczkach na $K$.

\emph{Odwrotność.}
Dla każdego wyczerpującego ciągu zbiorów
$A_l=U\cap B(0,l)$ i każdego $l$ możemy wybrać
zbiory zwarte $K_{l,j}\subset A_l$, na których
$f=F_{l,j}\in C^1$, z
$\lambda_n(A_l\setminus K_{l,j})<1/j$.
Prawie każdy punkt $A_l$ należy do przynajmniej
jednego z tych zbiorów i jest jego punktem gęstości.
W takim punkcie $f$ ma przybliżoną różniczkę
$DF_{l,j}$, ponieważ obie funkcje zgadzają się na
zbiorze gęstości $1$. Biorąc sumę po $l$,
otrzymujemy przybliżoną różniczkowalność prawie wszędzie.
\end{proof}

\begin{wniosek}[Luzinowskie wygładzenie odwzorowania Lipschitzowskiego]
\label{cor:dzety-lusin-lipschitz}
Niech $U\subset\mathbb R^n$ będzie otwarty,
$f:U\to\mathbb R^m$ lokalnie Lipschitzowskie,
a $A\subset U$ mierzalny o skończonej mierze.
Dla każdego $\varepsilon>0$ istnieje
$F\in C^1(\mathbb R^n,\mathbb R^m)$ takie, że
\[
 \lambda_n(\{x\in A:F(x)\ne f(x)\})<\varepsilon,
\]
a na pewnym zwartym $K\subset A$ o
$\lambda_n(A\setminus K)<\varepsilon$ zachodzi
$DF=Df$.
\end{wniosek}
\begin{proof}
Twierdzenie Rademachera daje przybliżoną
różniczkowalność prawie wszędzie.
Stosujemy Twierdzenie~\ref{thm:dzety-whitney-lusin}.
\end{proof}

\begin{uwaga}[Granice twierdzenia]
Twierdzenie~\ref{thm:dzety-whitney-lusin} nie mówi,
że $f$ jest funkcją $C^1$ po zmianie na zbiorze miary zero.
Wolno odrzucić zbiór o \emph{dowolnie małej}, lecz niekoniecznie
zerowej mierze, zależny od $\varepsilon$.
Na przykład $f(x)=|x|$ na $(-1,1)$ nie ma globalnego
przedstawiciela $C^1$ równego $f$ prawie wszędzie.
\end{uwaga}

\section{Dżety wyższych rzędów i zgodność danych}
Dla $k\ge2$ sam formalny zapis danych
$(a_\alpha(x))_{|\alpha|\le k}$ nie gwarantuje,
że są one pochodnymi jednej funkcji. W teorii miary
istotne jest, czy można znaleźć duży zbiór, na którym
dane spełniają \emph{wszystkie} warunki Whitneya.

\begin{twierdzenie}[Luzinowskie przedłużenie zgodnych dżetów rzędu $k$]
\label{thm:dzety-lusin-higher}
Niech $U\subset\mathbb R^n$ będzie otwarty, $k\ge1$,
a $a_\alpha:U\to\mathbb R$, $|\alpha|\le k$,
będą funkcjami mierzalnymi. Załóżmy, że dla
$\lambda_n$-prawie każdego $x\in U$ i każdego
$|\alpha|\le k$ zachodzi zwykła, względna do
pełnomiarowego zbioru $D\subset U$, ekspansja
\[
 a_\alpha(y)=\sum_{|\beta|\le k-|\alpha|}
 \frac{a_{\alpha+\beta}(x)}{\beta!}(y-x)^\beta
 +o(|y-x|^{k-|\alpha|})\quad(D\ni y\to x).
\]
Wówczas dla każdego mierzalnego $A\subset U$
o skończonej mierze i każdego $\varepsilon>0$
istnieją $K\subset A$ zwarty i
$F\in C^k(\mathbb R^n)$ takie, że
$\lambda_n(A\setminus K)<\varepsilon$ oraz
$D^\alpha F|_K=a_\alpha|_K$ dla wszystkich
$|\alpha|\le k$.
\end{twierdzenie}
\begin{proof}
Twierdzenie Luzina zastosowane do skończonej rodziny
$a_\alpha$ daje zwarty $K_0\subset A\cap D$ o dowolnie
małym dopełnieniu i ciągłości wszystkich współczynników.
Dla każdej wielopochodnej $\alpha$ i promienia $r$
weźmy supremum po $y\in K_0$ z $0<|x-y|<r$
ilorazu reszty Taylora przez $|x-y|^{k-|\alpha|}$.
Jest to funkcja mierzalna względem $x\in K_0$:
możemy zastąpić wszystkie $y$ przeliczalnym zbiorem
gęstym w $K_0$, ponieważ na $K_0$ współczynniki
są ciągłe. Z hipotezy supremum dąży do zera punktowo przy
$r\downarrow0$, ponieważ $K_0\subset D$ i każda
reszta jest oceniana wzdłuż całego zbioru $D$.
To właśnie tu korzystamy z silniejszej, względnej
wersji założenia, a nie jedynie ze zbieżności przybliżonej. Twierdzenie Jegorowa
pozwala przejść do zwartego $K\subset K_0$ z
arbitralnie małą utratą miary, na którym zbieżność
jest jednostajna dla skończenie wielu $\alpha$.
Właśnie to są warunki pola Whitneya z
Definicji~\ref{def:dzety-whitney-field}. Zastosowanie
Twierdzenia~\ref{thm:dzety-whitney-extension}
kończy dowód.
\end{proof}

\begin{uwaga}[Ważne rozróżnienie]
Warunek z Twierdzenia~\ref{thm:dzety-lusin-higher}
jest mocniejszy od samego istnienia przybliżonego
wielomianu Peana rzędu $k$ dla jednej funkcji.
Dla ogólnego twierdzenia o przybliżonych dżetach
wyższych rzędów potrzebny jest dodatkowy krok
przechodzący od informacji w mierze do zgodności
wszystkich pochodnych na wybranym zbiorze.
Nie wolno pomijać go, stosując twierdzenie Whitneya.
\end{uwaga}

\section{Dodatkowe przykłady i powiązania}
\begin{przyklad}[Dżet holomorficzny]
Niech $U\subset\mathbb C^n$ będzie obszarem,
$p\in U$ i $f\in\mathcal O(U)$. Jego holomorficzny
$k$-dżet zapisujemy w lokalnych współrzędnych jako
\[
 j_p^kf\equiv\sum_{|\alpha|\le k}
 \frac{\partial_z^\alpha f(p)}{\alpha!}(z-p)^\alpha
 \pmod{\mathfrak m_p^{k+1}}.
\]
Warunki interpolacji wartości i pochodnych
(np. w zadaniach typu Carathéodory'ego--Fejéra)
są warunkami na holomorficzne dżety.
W przeciwieństwie do dowolnych gładkich funkcji,
holomorficzne dżety podlegają globalnym ograniczeniom
wynikającym z istnienia funkcji holomorficznej na
ustalonym obszarze. Twierdzenie Whitneya w klasie $C^k$
nie zapewnia holomorficznego przedłużenia.
\end{przyklad}

\begin{przyklad}[Dżety a metoda elementów skończonych]
Przybliżenia typu Hermite'a w jednym wymiarze
łączą zadane wartości $u(x_i)$ i pochodne $u'(x_i)$
na końcach przedziałów. Są więc interpolacją
$1$-dżetów na skończonym zbiorze węzłów.
Warunki zgodności na węzłach decydują,
czy powstaje aproksymacja klasy $C^1$.
\end{przyklad}

\begin{przyklad}[Dżety i geometryczne oszacowania rzędu]
Dwie gładkie funkcje o identycznym $k$-dżecie w $p$
mogą różnić się globalnie, ale ich różnica
spełnia $f(x)-g(x)=o(|x-p|^k)$ przy $x\to p$.
Jeżeli mają również równe pochodne rzędu $k+1$,
to ich $(k+1)$-dżety są równe i otrzymujemy
odpowiednią poprawę lokalnego rzędu reszty.
\end{przyklad}

\section{Mapa zależności i zakres udowodnionych rezultatów}
\begin{center}
\begin{tikzpicture}[>=stealth, every node/.style={font=\small,align=center},
 box/.style={draw,rounded corners=2pt,inner sep=5pt,text width=3.8cm}]
 \node[box] (j) at (-2.4,2.8) {Dżety gładkie\\wiązki i formy kontaktowe};
 \node[box] (a) at (2.4,2.8) {Gęstość Lebesgue'a\\dżety przybliżone};
 \node[box] (w) at (-2.4,0.7) {Warunki Whitneya\\przedłużenie $C^k$};
 \node[box] (l) at (2.4,0.7) {Whitney--Luzin\\wybór dużego $K$};
 \node[box] (c) at (0,-1.45) {Zastosowania: Sobolew, $BV$,\\area/coarea, PDE};
 \draw[->] (j) -- (w);
 \draw[->] (a) -- (l);
 \draw[->] (w) -- (l);
 \draw[->] (l) -- (c);
\end{tikzpicture}
\end{center}

\begin{uwaga}[Co zostało udowodnione, a co przyjęte?]
W tym rozdziale szczegółowo wyprowadzono algebraiczny opis
dżetów, strukturę afiniczną, kryterium holonomiczności,
jednoznaczność dżetów przybliżonych, twierdzenie Whitneya
metodą kostek i twierdzenie Whitneya--Luzina dla $C^1$.
Osobno przyjęto znane twierdzenia Rademachera, Luzina,
Jegorowa i Sobolewa o przybliżonej różniczkowalności
oraz ogólne wzory area/coarea i rozkład $BV$.
Ich przytoczenie pokazuje zastosowania, nie stanowi
nowego dowodu tych wyników.
\end{uwaga}


\section*{Wskazówki bibliograficzne}
Konstrukcja wiązek dżetów i form kontaktowych jest standardowo
przedstawiana w monografii D.~J.~Saundersa,
\emph{The Geometry of Jet Bundles}. Dla teorii przedłużeń
podstawową pracą jest H.~Whitney,
\emph{Analytic extensions of differentiable functions defined
in closed sets} (1934). Przybliżona różniczkowalność,
twierdzenia typu Luzina i ich związek z teorią Sobolewa i $BV$
są omawiane w książce L.~C.~Evansa i R.~F.~Gariepy'ego,
\emph{Measure Theory and Fine Properties of Functions}.
W tekście rozdzielono twierdzenia dowodzone tutaj od
rezultatów zewnętrznych (Rademacher, własności Sobolewa,
wzory area i koobszaru). W konstrukcji Whitneya wykorzystuje
się standardowy rozkład na kostki i rozkład jedności;
podano ich własności oraz argument sklejania na zbiorze
domkniętym. Dowody powinny być czytane z zachowaniem
wskazanych założeń regularności.
